\begin{abstract}
Few-shot industrial anomaly detectors built on frozen foundation features often provide strong anomaly rankings, but their thresholds may not maintain consistent false-alarm rates across categories or input shifts. To address this gap, we introduce an auditable reliability framework with two complementary strengths: it exposes the exact resolution limit of target-only alarms and tests whether source-category evidence is sufficient for transferable certification rather than assuming that source pooling provides control. The framework combines a low-storage DINOv2 PCA residual with label-free leave-one-image-out (LOIO) p-values from $k$ normal supports. Evaluating 15 MVTec and 12 VisA categories under four corruption types, our audit reveals a resolution floor: target-only alarms cannot trigger below $\alpha=1/(k{+}1)$. Empirical normal p-values range from conservative on VisA at $k=4$ to anti-conservative on corrupted MVTec. We further introduce Cross-category Reliability Estimation with Source Support (CRESS), a source-assisted procedure that separates reference, threshold-proposal, and certification categories. Across a frozen grid of four $k$ values, five seeds, five conditions, within-dataset tests, and MVTec-to-VisA/MPDD transfer, strict category-level certification returns the fail-closed zero threshold in all 960 gate cells. This outcome is not solely attributable to Hoeffding looseness: when all observed category losses are zero, even an optimistic multiplicity-free, deterministic, uniformly valid, distribution-free 95\% upper bound requires at least 14 independent categories to certify $\alpha=0.20$. The frozen protocol requires more, whereas only three or four certification categories are available. Fixed-archive image certificates select positive thresholds in a fraction of settings ranging from 37\% to 60\%, but estimate image-level risk under the observed source mixture rather than new-category risk. Historical non-independent CRESS results are therefore retained only as empirical diagnostics. These results establish a clear feasibility boundary: the framework separates empirically useful source thresholds from transferable certificates and prevents unsupported claims of target-category control.
\end{abstract}
